% !PRINTED 175
\weberpage{175}
%<<CONT
variant $J$, by Theorem 4., as a rational function of $\Theta$,
then this function cannot change when one of the substitutions
$(\Theta, \Theta_1), (\Theta, \Theta_2), \dots$ is carried out, and this function
is therefore rationally expressible in terms of the coefficients of $\Phi(t)$.
Whether this representation is indeed possible by integral functions
would remain undecided in this proof.

Among the functions $\omega$ of the variables $x$ the
variables $x$ themselves are also contained, and the question from which we
set out, whether the variables $x$ can be regarded as algebraic functions
of the invariants, is thereby decided in the
affirmative.

The problem of representing the variables $x$ as algebraic functions
of the invariants of the group $S$, that is, the determination
of the field $\Omega(\Theta)$, is called, after \gesperrt{F. Klein}, the \gesperrt{form problem
of the group $S$}\footnote{Vorlesungen über das Ikosaëder, S.\,123. (Leipzig 1884.)}. If the number of variables is $n$, we call
it the form problem of the $n$th dimension.

The field $\Omega(\Theta)$ is an algebraic field over $\Omega$ completely determined
by the group $S$. It is a normal field,
for by 4. the conjugate quantities $\Theta, \Theta_1, \Theta_2, \dots \Theta_{\mu-1}$
are all contained in the field $\Omega(\Theta)$ itself. The equation $\Phi(t) = 0$
is a normal equation and is the \gesperrt{Galois resolvent
of the form problem} (Bd.\ I, \S.\,145).

\gesperrt{5. The Galois group of the form problem,
i.\,e.\ the Galois group of the equation
$\Phi(t) = 0$, is isomorphic with the group $S$.}

To prove this, we denote by $A, B$ two substitutions
from $S$ and by $AB = C$ the substitution composed
from them. Now if
\[ \Theta_1 = \Theta[A(x)],\ \Theta_2 = \Theta[B(x)],\ \Theta_3 = \Theta[C(x)], \]
then the substitution
\[ (\Theta, \Theta_2) = (\Theta_1, \Theta_3), \]
and consequently
\[ (\Theta, \Theta_1)(\Theta, \Theta_2) = (\Theta, \Theta_1)(\Theta_1, \Theta_3) = (\Theta, \Theta_3), \]
i.\,e.\ the group of the substitutions $(\Theta, \Theta_1), (\Theta, \Theta_2), \dots$ is
isomorphic with the group of the $A, B, \dots$.

If instead of the function $\Theta$ we take a function $\eta$ which does not
have all its values distinct, but admits the substitutions of a
%CONT>>
